\documentclass[../main]{subfiles}
\begin{document}
\chapter{Conceptos básicos de dibujo 3D}
\includegraphics[scale=0.35]{./images/05/ejemplo.jpg}
\section{Extrusión y Redondeado}

\subsection{Llavero}
 Dibujar un llavero de $60mm \times 25mm \times 6mm$, dentro debe tener su nombre en forma de relieve $h=2mm$, y debe estar perforado con un diámetro de $d=5mm$ en un extremo para poder colgar llaves. Los cantos deben estar redondeados $r=0,5cm$.

\includegraphics[]{./images/05/llavero.png}



    
\subsection{Taza con asa}

Dibujar una taza con asa, el díametro de la taza es de $d_{e}=8cm$ y la altura $h=9cm$ el diámetro interior es de $d_{i}=7cm$, debajo de la taza debe escribir su nombre y apellido en $h=1mm$ de profundidad de hendidura. El asa o \emph{la oreja} debe ser realizado con un toroide de $5mm de radio \times 2,5cm$ para mayor comodidad de uso.

\includegraphics[]{./images/05/tazaasa.png}


\subsection{Dado personal} Realizar un dado de $3cm\times 3cm \times 3cm$, con cantos redondeados a $r=5mm$ y letras de altura $h_{letra}=2cm$, $h_{profundidad}=2mm$. En vez de colocar números coloque las primeras 6 letras de su apellido.
    \subsubsection{Pasos:}
    \begin{enumerate}
        \item Realice un cubo de $3cm\times 3cm \times 3cm$
        \item Redondar vértices del cubo.
        \item Coloque los letas en cada cara, colocando la forma a partir de texto en cada cara.
        \item Extruir la forma de los textos.
        \item Restar la los textos extruidos al cubo.
    \end{enumerate}
 
 \subsection{Soporte para celular}
 Realizar un soporte \emph{a medida de su celular}, bajo el siguiente concepto. En la base escribir su nombre y apellido.
\img{images/05/soporte.png}{Soporte para celular}
 
\subsection{Dibujo de portada}
Dibujar la imagen de la portada en del práctico como pieza en 3D, respetando sus medidas acotadas. 

\subsection{Dibujar las siguientes piezas}
\begin{center}
\includegraphics[scale=0.5]{./images/05/01.jpg}

\includegraphics[scale=1]{./images/05/04.png}

\includegraphics[scale=0.8]{./images/05/06.png}

\includegraphics[scale=1]{./images/05/02.jpg}

\includegraphics[scale=1]{./images/05/03.jpg}

\includegraphics[scale=0.7]{./images/05/05.jpg}

\includegraphics[scale=1.4]{./images/05/07.jpg}

\end{center}

\end{document}